% ART_RECOMB: additional figures. Upload the fig_*.pdf files to a figures/ folder in Overleaf.
% Packages already in the preamble: graphicx. Paste this after Appendix D.
% The body is at the 10-page limit, so everything here belongs in the appendix.
% To use the method overview in the body instead, see the note in the hand-off message and
% recompile: going past 10 body pages is reject-without-review.

\section{Additional figures}

\begin{figure}[p]
\centering
\includegraphics[width=\textwidth]{figures/fig_method.pdf}
\caption{The selection-aware test and the evidence gate. The locus $L$ and every shuffle under
$\mathcal{M}_{\mathrm{dinuc}}$ and $\mathcal{M}_{\mathrm{block}}$ pass through the same procedure $D$, so the
reference distribution is that of the search outcome. The larger of the two empirical $p$-values is
corrected over a declared family, and the gate keeps array evidence in its own field. Histograms are
the 1,000 draws per background for LPJP1 QXN67953.1.}
\label{fig:method}
\end{figure}

\begin{figure}[p]
\centering
\includegraphics[width=\textwidth]{figures/fig_calibration.pdf}
\caption{Behaviour of the null. \textbf{a}, Observed best-chain score against the pooled full-search null
for each benchmark family and reference locus (6\,kb window, 1,000 draws per background). The filled
diamond for MarsHill is the chain reported as its array; the hollow diamond is its top-scoring chain.
\textbf{b}, On the DGR control the local block background alone gives $\hat p=0.003$ and the dinucleotide
background $\hat p=0.066$; carrying the larger withholds the call. \textbf{c}, Each of the 60,000 stored
draws was treated as a pseudo-observed score and tested against its own background (leaving it out),
against the other background, and against both by the maximum rule, across all 30 windows. At
$\alpha=0.001$ the other background rejects at 1.5 times nominal on average and at 10 times nominal at the
97.5th percentile of window--background combinations (shaded). The maximum cannot exceed the
own-background rate, so it stays at or below $\alpha$. These draws come from the stated models, so this
measures sensitivity to the choice of background, not validity on genomic sequence.
\textbf{d}, Smallest attainable corrected value by draw count, with the declared families used here.
\textbf{e}, Corrected values for the seven published references as the family grows: 5 of 7 reach 0.05 at
$m=9$ and 2 of 7 at $m=31$. Solid lines use each top-scoring chain; the dashed line is MarsHill's reported
chain.}
\label{fig:calibration}
\end{figure}

\begin{figure}[p]
\centering
\includegraphics[width=\textwidth]{figures/fig_benchmark.pdf}
\caption{All 30 designed cases, by position in the assessed window. Planted copies are drawn above each line
and detected copies below it; red triangles mark detections matching no planted copy. Columns give copy
counts, corrected $p$ ($m=30$, floor $30/1001\approx0.030$) and whether the expected behaviour was met. Each
of the three failures recovers every planted copy and adds one unmatched detection.
\texttt{repetitive\_DNA\_2} reports no copies despite a floor-level search score; it passes because its
expected behaviour is to withhold an ART call.}
\label{fig:benchmark}
\end{figure}

% RECONCILE BEFORE SUBMITTING: Table 1 reports tuned PILER-CR calling 24/31 inputs; the wrapper run in
% panel b calls 22/31. The wrapper's own REPORT.md flags this and the original harness was not available.
\begin{figure}[p]
\centering
\includegraphics[width=\textwidth]{figures/fig_baselines.pdf}
\caption{Search and calibration are separate comparisons. \textbf{a}, Copy recall and precision from
Table~\ref{tab:baselines}. \textbf{b}, Each of the 31 inputs with copy-level truth, tested under one wrapper
with the same 1,000 shuffles per background. The PILER-CR statistic is $(\text{copies}-1)\times$ consensus
repeat length $\times$ mean copy identity of its best array, and 0 if it reports none. Wrapped PILER-CR is
close to $D$ on designed cases, missing \texttt{missing\_copy\_1} and \texttt{\_2}, and has no power on the
references because it calls no array there. It segfaulted on 14 of the 62,000 shuffles, all from reference
windows; counting these as exceedances or as no array gives identical counts. Bonferroni is over all 31
inputs, with floor $31/1001=0.031$.}
\label{fig:baselines}
\end{figure}

\begin{figure}[p]
\centering
\includegraphics[width=\textwidth]{figures/fig_references.pdf}
\caption{The nine reference RT loci. \textbf{a}, Oriented loci aligned at the RT start, supported above the
dashed line and held for review below it, with $\tilde p$ corrected over $m=9$. The reasons for holding are
visible in the annotation: Madawaska's repeats overlap annotated CDSs and LY01 has an annotated gene between
array and RT. \textbf{b}, Difference between our spacers and those implied by the 22-nt core-centred windows
of the source work, for all 28 spacers of the seven loci with published coordinates. The offset equals the
difference in repeat width between the two definitions, so core recovery does not establish boundary
reproduction.}
\label{fig:references}
\end{figure}

\begin{figure}[p]
\centering
\includegraphics[width=\textwidth]{figures/fig_geometry.pdf}
\caption{Array geometry and repeat content. \textbf{a}, Repeat width against the gap between consecutive
repeats, for the seven \emph{Staphylococcus} references and the four supported type I tracts, with vertical
jitter for visibility. The grey band is the default spacer range of the finders and the dashed box the tuned
range of Table~\ref{tab:baselines}. Gaps of 630 and 689\,nt fall outside the tuned range; the 689-nt gap is
flagged as possibly spanning missing repeats. \textbf{b}, Repeat logos against a uniform background, with no
small-sample correction, so they are descriptive at 3--5 copies. \emph{Staphylococcus} copies are stacked in
register on the published core, which is used for scoring only; columns covered by fewer than half of the
copies are omitted.}
\label{fig:geometry}
\end{figure}

\begin{figure}[p]
\centering
\includegraphics[width=\textwidth]{figures/fig_acquisition.pdf}
\caption{From bounded sources to the released dataset. The expansion adds 21 distinct type I pairs; with the 13
catalogued and 2 follow-up type I pairs and 10 type II pairs this gives 49 raw records and, after collapsing 3
aliases, 46 pairs. Unresolved retrievals are recorded as such and are not absences. The counts are not
prevalence estimates. The last two tiers restate the matching graph of Figure~2a.}
\label{fig:acquisition}
\end{figure}

\begin{figure}[p]
\centering
\includegraphics[width=\textwidth]{figures/fig_trees.pdf}
\caption{Sequence-only placement. \textbf{a}, RT-core tree of 137 sequences (MAFFT L-INS-i, 268 columns, FastTree
LG+$\Gamma$), midpoint-rooted for display. \textbf{b}, The RT-core tree pruned to the 36 type I pairs, against
their partner tree (626 columns), with nodes rotated to reduce crossings. Labels are the last two digits of the
ARTS identifier. Internal values are local SH-like support, not bootstraps. Placement is descriptive: it does not
establish independent lineages, co-evolution, or ART rather than NART membership.}
\label{fig:trees}
\end{figure}

\begin{figure}[p]
\centering
\includegraphics[width=\textwidth]{figures/fig_landscape.pdf}
\caption{The 46 pairs, by ARTS identifier (last two digits). \textbf{a}, Upstream DNA available in each locus,
coloured by recorded array status; statuses keep their original cohort assessments and were not retested on the
merged dataset. \textbf{b}, Protein lengths, with RT above the axis and partner below. The NTD is operational:
the residues before the first significant RVT\_1 envelope.}
\label{fig:landscape}
\end{figure}

% OPTIONAL, left out by default. Enabling this means the sentence "we claim no structural or interface
% result here" needs a clause saying a partial exploratory prediction run is shown in this appendix.
% Regenerate from the finished run first: python make_figures.py binding
\iffalse
\begin{figure}[p]
\centering
\includegraphics[width=\textwidth]{figures/fig_binding.pdf}
\caption{Partial, exploratory predicted-interface run (hosted Boltz-2; 56 of 1,221 requests and 168 of 3,663
models returned). \textbf{a}, Median ipTM for full-length type I pairs without alignments; native pairs are
outlined and missing cells are hatched. No row of this matrix supports a native-versus-alternative contrast and
none is computed. Returned medians span 0.12--0.25. \textbf{b}, \textbf{c}, Adding independently ordered per-chain
alignments raises per-residue confidence on the two pairs run both ways, while interface confidence stays low.
\textbf{d}, NTD share of 5\,\AA{} contacts against NTD share of RT residues for the 18 native query-only pairs,
with the line of equal shares. \textbf{e}, C$\alpha$ trace of one model, lightly smoothed for display. These are
predicted structures and confidence scores from stochastic samples: they do not measure binding, and low
confidence does not establish its absence.}
\label{fig:binding}
\end{figure}
\fi
